\documentclass[12pt,a4paper]{report}

% ============================================================================
% PAKIETY
% ============================================================================
\usepackage[utf8]{inputenc}
\usepackage[T1]{fontenc}
\usepackage[polish]{babel}
\usepackage{geometry}
\usepackage{graphicx}
\usepackage{listings}
\usepackage{xcolor}
\usepackage{hyperref}
\usepackage{booktabs}
\usepackage{longtable}
\usepackage{tabularx}
\usepackage{float}
\usepackage{caption}
\usepackage{subcaption}
\usepackage{fancyhdr}
\usepackage{titlesec}
\usepackage{tocloft}
\usepackage{enumitem}
\usepackage{amsmath}
\usepackage{tikz}
\usetikzlibrary{shapes.geometric, arrows, positioning, fit, calc}

% ============================================================================
% KONFIGURACJA
% ============================================================================
\geometry{margin=2.5cm}
\setlength{\parindent}{0pt}
\setlength{\parskip}{6pt}

% Kolory dla kodu
\definecolor{codegreen}{rgb}{0,0.6,0}
\definecolor{codegray}{rgb}{0.5,0.5,0.5}
\definecolor{codepurple}{rgb}{0.58,0,0.82}
\definecolor{backcolour}{rgb}{0.95,0.95,0.92}
\definecolor{javared}{rgb}{0.6,0,0}
\definecolor{javablue}{rgb}{0,0,0.6}

% Styl dla kodu Java
\lstdefinestyle{javastyle}{
    backgroundcolor=\color{backcolour},
    commentstyle=\color{codegreen},
    keywordstyle=\color{javablue}\bfseries,
    numberstyle=\tiny\color{codegray},
    stringstyle=\color{javared},
    basicstyle=\ttfamily\footnotesize,
    breakatwhitespace=false,
    breaklines=true,
    captionpos=b,
    keepspaces=true,
    numbers=left,
    numbersep=5pt,
    showspaces=false,
    showstringspaces=false,
    showtabs=false,
    tabsize=2,
    frame=single,
    language=Java
}

% Styl dla SQL
\lstdefinestyle{sqlstyle}{
    backgroundcolor=\color{backcolour},
    commentstyle=\color{codegreen},
    keywordstyle=\color{javablue}\bfseries,
    numberstyle=\tiny\color{codegray},
    stringstyle=\color{javared},
    basicstyle=\ttfamily\footnotesize,
    breakatwhitespace=false,
    breaklines=true,
    captionpos=b,
    keepspaces=true,
    numbers=left,
    numbersep=5pt,
    showspaces=false,
    showstringspaces=false,
    showtabs=false,
    tabsize=2,
    frame=single,
    language=SQL,
    morekeywords={SERIAL, BIGSERIAL, VARCHAR, BOOLEAN, TIMESTAMP, REFERENCES, CASCADE, RESTRICT}
}

% Styl dla properties
\lstdefinestyle{propstyle}{
    backgroundcolor=\color{backcolour},
    commentstyle=\color{codegreen},
    basicstyle=\ttfamily\footnotesize,
    breaklines=true,
    frame=single
}

\lstset{style=javastyle}

% Nagłówki i stopki
\pagestyle{fancy}
\fancyhf{}
\fancyhead[L]{\leftmark}
\fancyhead[R]{System Zarządzania Projektami}
\fancyfoot[C]{\thepage}

% Hyperref
\hypersetup{
    colorlinks=true,
    linkcolor=blue,
    filecolor=magenta,
    urlcolor=cyan,
    pdftitle={Dokumentacja Systemu Zarządzania Projektami},
    pdfauthor={Autor}
}

% ============================================================================
% STRONA TYTUŁOWA
% ============================================================================
\begin{document}

\begin{titlepage}
    \centering
    \vspace*{2cm}
    
    {\scshape\LARGE Politechnika Bydgoska \par}
    \vspace{0.5cm}
    {\scshape\Large Wydział Informatyki \par}
    \vspace{2cm}
    
    {\huge\bfseries System Zarządzania Projektami\par}
    \vspace{0.5cm}
    {\Large Aplikacja webowa REST API\par}
    \vspace{2cm}
    
    {\Large\bfseries Dokumentacja Projektu\par}
    \vspace{1cm}
    
    {\large Przedmiot: Bazy Danych\par}
    \vspace{2cm}
    
    \begin{tabular}{rl}
        \textbf{Autor:} & [Imię i Nazwisko] \\
        \textbf{Nr albumu:} & [Numer albumu] \\
        \textbf{Kierunek:} & Informatyka \\
        \textbf{Semestr:} & [Numer semestru] \\
    \end{tabular}
    
    \vfill
    
    {\large Bydgoszcz, \the\year\par}
\end{titlepage}

% ============================================================================
% SPIS TREŚCI
% ============================================================================
\tableofcontents
\newpage

% ============================================================================
% ROZDZIAŁ 1: OPIS SŁOWNY REALIZOWANEGO ZADANIA
% ============================================================================
\chapter{Opis słowny realizowanego zadania}

\section{Cel projektu}

Celem projektu jest stworzenie kompleksowego systemu zarządzania projektami (Project Management System), który umożliwia zespołom efektywną organizację pracy, śledzenie postępów oraz współpracę w czasie rzeczywistym.

System został zaimplementowany jako aplikacja typu REST API z wykorzystaniem nowoczesnych technologii backendowych, zapewniając:
\begin{itemize}
    \item Skalowalność i wydajność
    \item Bezpieczeństwo danych użytkowników
    \item Łatwą integrację z aplikacjami frontendowymi
    \item Pełną dokumentację API (OpenAPI/Swagger)
\end{itemize}

\section{Zakres funkcjonalności}

System realizuje następujące funkcjonalności:

\subsection{Moduł autentykacji i autoryzacji}
\begin{itemize}
    \item Rejestracja nowych użytkowników z walidacją danych
    \item Logowanie z wykorzystaniem tokenów JWT (JSON Web Token)
    \item Odświeżanie tokenów dostępu
    \item Resetowanie hasła przez email
    \item Role użytkowników: USER, ADMIN
\end{itemize}

\subsection{Moduł zarządzania projektami}
\begin{itemize}
    \item Tworzenie, edycja i usuwanie projektów
    \item Przypisywanie użytkowników do projektów z rolami (OWNER, MEMBER, VIEWER)
    \item Zarządzanie statusami projektów (NOT\_STARTED, IN\_PROGRESS, ON\_HOLD, COMPLETED, CANCELLED)
    \item Określanie dat rozpoczęcia i zakończenia projektu
\end{itemize}

\subsection{Moduł zarządzania zadaniami}
\begin{itemize}
    \item CRUD operacje na zadaniach w ramach projektów
    \item Przypisywanie zadań do użytkowników
    \item Statusy zadań: TODO, IN\_PROGRESS, REVIEW, DONE
    \item Terminy realizacji zadań
\end{itemize}

\subsection{Moduł komentarzy}
\begin{itemize}
    \item Komentarze do projektów
    \item Komentarze do zadań
    \item Edycja i usuwanie własnych komentarzy
\end{itemize}

\subsection{Moduł plików}
\begin{itemize}
    \item Upload plików do projektów
    \item Pobieranie plików
    \item Walidacja typów i rozmiarów plików
    \item Zarządzanie avatarami użytkowników
\end{itemize}

\subsection{Moduł powiadomień}
\begin{itemize}
    \item Automatyczne powiadomienia o przypisaniu do projektu/zadania
    \item Powiadomienia o zmianach statusów
    \item Oznaczanie powiadomień jako przeczytane
\end{itemize}

\subsection{Moduł audytu}
\begin{itemize}
    \item Automatyczne logowanie wszystkich operacji
    \item Historia zmian w systemie
    \item Statystyki aktywności użytkowników
\end{itemize}

\subsection{Moduł komunikacji}
\begin{itemize}
    \item Chat w czasie rzeczywistym (WebSocket)
    \item Historia wiadomości
\end{itemize}

\section{Opis problemu biznesowego}

W dzisiejszych czasach efektywne zarządzanie projektami stanowi kluczowy element sukcesu organizacji. Tradycyjne metody komunikacji i śledzenia postępów (email, arkusze kalkulacyjne) są niewystarczające dla zespołów rozproszonych geograficznie.

System rozwiązuje następujące problemy:
\begin{enumerate}
    \item \textbf{Brak centralnego miejsca informacji} -- wszystkie dane projektowe w jednym systemie
    \item \textbf{Trudności w śledzeniu postępów} -- statusy zadań i projektów w czasie rzeczywistym
    \item \textbf{Problemy z komunikacją} -- wbudowany chat i komentarze
    \item \textbf{Brak historii zmian} -- pełny audyt wszystkich operacji
    \item \textbf{Zarządzanie dokumentami} -- centralne repozytorium plików projektowych
\end{enumerate}

\section{Stos technologiczny}

\begin{table}[H]
\centering
\caption{Wykorzystane technologie}
\begin{tabular}{|l|l|l|}
\hline
\textbf{Warstwa} & \textbf{Technologia} & \textbf{Wersja} \\
\hline
Backend & Spring Boot & 3.4.4 \\
Język programowania & Java & 23 \\
Baza danych & PostgreSQL & 15+ \\
ORM & Spring Data JPA / Hibernate & 6.x \\
Autentykacja & JWT (jjwt) & 0.11.5 \\
Dokumentacja API & OpenAPI / Swagger & 2.5.0 \\
WebSocket & Spring WebSocket & -- \\
Email & Spring Mail + Thymeleaf & -- \\
Build tool & Gradle & 8.13 \\
Testy & JUnit 5 + TestContainers & -- \\
\hline
\end{tabular}
\end{table}

% ============================================================================
% ROZDZIAŁ 2: OPIS FORMALNY
% ============================================================================
\chapter{Opis formalny systemu}

\section{Diagram przypadków użycia}

Poniższy diagram przedstawia główne przypadki użycia systemu z podziałem na aktorów.

\begin{figure}[H]
\centering
\begin{tikzpicture}[
    actor/.style={circle, draw, minimum size=1cm},
    usecase/.style={ellipse, draw, minimum width=2.5cm, minimum height=1cm, text width=2.3cm, align=center, font=\scriptsize},
    system/.style={rectangle, draw, dashed, minimum width=10cm, minimum height=12cm}
]

% System boundary
\node[system, label=above:{\textbf{System Zarządzania Projektami}}] (sys) at (5,0) {};

% Actors
\node[actor, label=below:{Gość}] (guest) at (-1,2) {G};
\node[actor, label=below:{Użytkownik}] (user) at (-1,-2) {U};
\node[actor, label=below:{Admin}] (admin) at (11,-2) {A};

% Use cases - Authentication
\node[usecase] (uc1) at (3,4) {Rejestracja};
\node[usecase] (uc2) at (3,2.5) {Logowanie};
\node[usecase] (uc3) at (3,1) {Reset hasła};

% Use cases - Projects
\node[usecase] (uc4) at (3,-0.5) {Zarządzanie projektami};
\node[usecase] (uc5) at (3,-2) {Zarządzanie zadaniami};

% Use cases - Other
\node[usecase] (uc6) at (7,4) {Komentarze};
\node[usecase] (uc7) at (7,2.5) {Pliki};
\node[usecase] (uc8) at (7,1) {Powiadomienia};
\node[usecase] (uc9) at (7,-0.5) {Chat};
\node[usecase] (uc10) at (7,-2) {Audyt};

% Connections - Guest
\draw (guest) -- (uc1);
\draw (guest) -- (uc2);
\draw (guest) -- (uc3);

% Connections - User
\draw (user) -- (uc4);
\draw (user) -- (uc5);
\draw (user) -- (uc6);
\draw (user) -- (uc7);
\draw (user) -- (uc8);
\draw (user) -- (uc9);

% Connections - Admin
\draw (admin) -- (uc10);
\draw (admin) -- (uc4);

% Inheritance
\draw[dashed, ->] (user) -- (guest);
\draw[dashed, ->] (admin) -- (user);

\end{tikzpicture}
\caption{Diagram przypadków użycia}
\end{figure}

\section{Diagram ERD (Entity-Relationship Diagram)}

\begin{figure}[H]
\centering
\begin{tikzpicture}[
    entity/.style={rectangle, draw, minimum width=3cm, minimum height=0.8cm, fill=blue!10},
    attribute/.style={ellipse, draw, minimum width=1.5cm, fill=yellow!20, font=\tiny},
    relation/.style={diamond, draw, minimum width=1.5cm, fill=green!20, font=\scriptsize}
]

% Entities
\node[entity] (users) at (0,0) {\textbf{USERS}};
\node[entity] (projects) at (6,0) {\textbf{PROJECTS}};
\node[entity] (tasks) at (6,-4) {\textbf{TASKS}};
\node[entity] (comments) at (0,-4) {\textbf{COMMENTS}};
\node[entity] (files) at (10,-2) {\textbf{FILES}};
\node[entity] (notifications) at (-4,-2) {\textbf{NOTIFICATIONS}};
\node[entity] (audit) at (3,-6) {\textbf{AUDIT\_LOGS}};

% Relations
\node[relation] (rel1) at (3,0) {tworzy};
\node[relation] (rel2) at (6,-2) {zawiera};
\node[relation] (rel3) at (3,-4) {ma};
\node[relation] (rel4) at (8,-0.5) {posiada};
\node[relation] (rel5) at (-2,-1) {otrzymuje};

% Connections
\draw (users) -- (rel1) node[midway,above] {1};
\draw (rel1) -- (projects) node[midway,above] {N};
\draw (projects) -- (rel2) node[midway,right] {1};
\draw (rel2) -- (tasks) node[midway,right] {N};
\draw (tasks) -- (rel3) node[midway,above] {1};
\draw (rel3) -- (comments) node[midway,above] {N};
\draw (projects) -- (rel4) node[midway,above] {1};
\draw (rel4) -- (files) node[midway,above] {N};
\draw (users) -- (rel5) node[midway,left] {1};
\draw (rel5) -- (notifications) node[midway,left] {N};

% User-Audit
\draw (users) -- (audit) node[midway,left] {1:N};

\end{tikzpicture}
\caption{Uproszczony diagram ERD}
\end{figure}

\section{Definicje tabel i typy danych}

\subsection{Tabela USERS}

\begin{table}[H]
\centering
\caption{Struktura tabeli USERS}
\begin{tabular}{|l|l|l|l|l|}
\hline
\textbf{Kolumna} & \textbf{Typ} & \textbf{Null} & \textbf{Klucz} & \textbf{Opis} \\
\hline
id & BIGSERIAL & NOT NULL & PK & Identyfikator \\
first\_name & VARCHAR(50) & NOT NULL & -- & Imię \\
last\_name & VARCHAR(50) & NOT NULL & -- & Nazwisko \\
email & VARCHAR(100) & NOT NULL & UQ & Email (unikalny) \\
password & VARCHAR(255) & NOT NULL & -- & Hash hasła (BCrypt) \\
role & VARCHAR(20) & NOT NULL & -- & Rola (USER/ADMIN) \\
avatar\_file\_name & VARCHAR(255) & NULL & -- & Nazwa pliku avatara \\
avatar\_content\_type & VARCHAR(50) & NULL & -- & Typ MIME avatara \\
created\_at & TIMESTAMP & NOT NULL & -- & Data utworzenia \\
\hline
\end{tabular}
\end{table}

\subsection{Tabela PROJECTS}

\begin{table}[H]
\centering
\caption{Struktura tabeli PROJECTS}
\begin{tabular}{|l|l|l|l|l|}
\hline
\textbf{Kolumna} & \textbf{Typ} & \textbf{Null} & \textbf{Klucz} & \textbf{Opis} \\
\hline
id & BIGSERIAL & NOT NULL & PK & Identyfikator \\
name & VARCHAR(255) & NOT NULL & -- & Nazwa projektu \\
description & VARCHAR(1000) & NULL & -- & Opis projektu \\
start\_date & DATE & NULL & -- & Data rozpoczęcia \\
end\_date & DATE & NULL & -- & Data zakończenia \\
status & VARCHAR(20) & NOT NULL & -- & Status projektu \\
created\_by & BIGINT & NOT NULL & FK & ID twórcy \\
created\_at & DATE & NOT NULL & -- & Data utworzenia \\
\hline
\end{tabular}
\end{table}

\subsection{Tabela TASKS}

\begin{table}[H]
\centering
\caption{Struktura tabeli TASKS}
\begin{tabular}{|l|l|l|l|l|}
\hline
\textbf{Kolumna} & \textbf{Typ} & \textbf{Null} & \textbf{Klucz} & \textbf{Opis} \\
\hline
id & BIGSERIAL & NOT NULL & PK & Identyfikator \\
name & VARCHAR(100) & NOT NULL & -- & Nazwa zadania \\
description & VARCHAR(500) & NULL & -- & Opis zadania \\
status & VARCHAR(20) & NOT NULL & -- & Status zadania \\
due\_date & DATE & NULL & -- & Termin realizacji \\
project\_id & BIGINT & NOT NULL & FK & ID projektu \\
assigned\_to & BIGINT & NULL & FK & ID przypisanego użytk. \\
created\_at & DATE & NOT NULL & -- & Data utworzenia \\
\hline
\end{tabular}
\end{table}

\subsection{Tabela PROJECT\_USER (relacja N:M)}

\begin{table}[H]
\centering
\caption{Struktura tabeli PROJECT\_USER}
\begin{tabular}{|l|l|l|l|l|}
\hline
\textbf{Kolumna} & \textbf{Typ} & \textbf{Null} & \textbf{Klucz} & \textbf{Opis} \\
\hline
project\_id & BIGINT & NOT NULL & PK, FK & ID projektu \\
user\_id & BIGINT & NOT NULL & PK, FK & ID użytkownika \\
role & VARCHAR(20) & NOT NULL & -- & Rola w projekcie \\
\hline
\end{tabular}
\end{table}

\subsection{Tabela AUDIT\_LOGS}

\begin{table}[H]
\centering
\caption{Struktura tabeli AUDIT\_LOGS}
\begin{tabular}{|l|l|l|l|l|}
\hline
\textbf{Kolumna} & \textbf{Typ} & \textbf{Null} & \textbf{Klucz} & \textbf{Opis} \\
\hline
id & BIGSERIAL & NOT NULL & PK & Identyfikator \\
user\_id & BIGINT & NOT NULL & FK & ID użytkownika \\
action & VARCHAR(50) & NOT NULL & -- & Typ akcji \\
entity & VARCHAR(100) & NOT NULL & -- & Nazwa encji \\
entity\_id & BIGINT & NOT NULL & -- & ID encji \\
details & VARCHAR(1000) & NULL & -- & Szczegóły \\
ip\_address & VARCHAR(45) & NULL & -- & Adres IP \\
timestamp & TIMESTAMP & NOT NULL & -- & Czas operacji \\
\hline
\end{tabular}
\end{table}

\subsection{Tabela NOTIFICATIONS}

\begin{table}[H]
\centering
\caption{Struktura tabeli NOTIFICATIONS}
\begin{tabular}{|l|l|l|l|l|}
\hline
\textbf{Kolumna} & \textbf{Typ} & \textbf{Null} & \textbf{Klucz} & \textbf{Opis} \\
\hline
id & BIGSERIAL & NOT NULL & PK & Identyfikator \\
user\_id & BIGINT & NOT NULL & FK & ID użytkownika \\
message & VARCHAR(500) & NOT NULL & -- & Treść powiadomienia \\
read & BOOLEAN & NOT NULL & -- & Czy przeczytane \\
type & VARCHAR(50) & NULL & -- & Typ powiadomienia \\
related\_entity\_id & BIGINT & NULL & -- & ID powiązanej encji \\
related\_entity\_type & VARCHAR(100) & NULL & -- & Typ powiązanej encji \\
created\_at & DATE & NOT NULL & -- & Data utworzenia \\
\hline
\end{tabular}
\end{table}

\section{Główne kwerendy systemu}

\subsection{Kwerendy JPQL w repozytoriach}

\begin{lstlisting}[style=javastyle, caption={ProjectUserRepository - zapytania}]
// Pobranie uzytkownikow projektu z eagerly loaded User
@Query("SELECT pu FROM ProjectUser pu 
        JOIN FETCH pu.user 
        WHERE pu.project.id = :projectId")
List<ProjectUser> findByProjectIdWithUsers(
    @Param("projectId") Long projectId);

// Sprawdzenie czy uzytkownik nalezy do projektu
boolean existsByProjectIdAndUserId(
    Long projectId, Long userId);
\end{lstlisting}

\begin{lstlisting}[style=javastyle, caption={AuditLogRepository - zapytania}]
// Pobranie logow z zakresu dat
@Query("SELECT a FROM AuditLog a 
        WHERE a.timestamp BETWEEN :startDate AND :endDate 
        ORDER BY a.timestamp DESC")
List<AuditLog> findByTimestampBetween(
    @Param("startDate") LocalDateTime startDate, 
    @Param("endDate") LocalDateTime endDate);

// Pobranie ostatnich logow uzytkownika
@Query("SELECT a FROM AuditLog a 
        WHERE a.user.id = :userId AND a.timestamp >= :since 
        ORDER BY a.timestamp DESC")
List<AuditLog> findRecentByUserId(
    @Param("userId") Long userId, 
    @Param("since") LocalDateTime since);

// Zliczanie akcji dla encji
@Query("SELECT COUNT(a) FROM AuditLog a 
        WHERE a.entity = :entity AND a.action = :action")
Long countByEntityAndAction(
    @Param("entity") String entity, 
    @Param("action") String action);
\end{lstlisting}

\begin{lstlisting}[style=javastyle, caption={NotificationRepository - zapytania}]
// Oznaczenie wszystkich jako przeczytane
@Modifying
@Query("UPDATE Notification n 
        SET n.read = true 
        WHERE n.user.id = :userId AND n.read = false")
int markAllAsReadByUserId(@Param("userId") Long userId);

// Zliczanie nieprzeczytanych
@Query("SELECT COUNT(n) FROM Notification n 
        WHERE n.user.id = :userId AND n.read = false")
Long countUnreadByUserId(@Param("userId") Long userId);
\end{lstlisting}

\subsection{Widoki SQL (Views)}

\begin{lstlisting}[style=sqlstyle, caption={Widok statystyk projektów}]
CREATE OR REPLACE VIEW v_project_statistics AS
SELECT 
    p.id AS project_id,
    p.name AS project_name,
    p.status,
    COUNT(DISTINCT t.id) AS total_tasks,
    COUNT(DISTINCT CASE WHEN t.status = 'DONE' 
          THEN t.id END) AS completed_tasks,
    COUNT(DISTINCT pu.user_id) AS team_members,
    COUNT(DISTINCT pc.id) AS total_comments,
    COUNT(DISTINCT pf.id) AS total_files
FROM projects p
LEFT JOIN tasks t ON t.project_id = p.id
LEFT JOIN project_user pu ON pu.project_id = p.id
LEFT JOIN project_comments pc ON pc.project_id = p.id
LEFT JOIN project_files pf ON pf.project_id = p.id
GROUP BY p.id, p.name, p.status;
\end{lstlisting}

\begin{lstlisting}[style=sqlstyle, caption={Widok zadań po terminie}]
CREATE OR REPLACE VIEW v_overdue_tasks AS
SELECT 
    t.id AS task_id,
    t.name AS task_name,
    t.due_date,
    t.status,
    p.name AS project_name,
    CONCAT(u.first_name, ' ', u.last_name) AS assigned_to_name,
    CURRENT_DATE - t.due_date AS days_overdue
FROM tasks t
JOIN projects p ON p.id = t.project_id
LEFT JOIN users u ON u.id = t.assigned_to
WHERE t.due_date < CURRENT_DATE 
  AND t.status != 'DONE'
ORDER BY days_overdue DESC;
\end{lstlisting}

\section{Formularze i endpointy API}

\subsection{Endpointy REST API}

\begin{longtable}{|l|l|p{6cm}|}
\caption{Lista głównych endpointów API} \\
\hline
\textbf{Metoda} & \textbf{Endpoint} & \textbf{Opis} \\
\hline
\endfirsthead
\hline
\textbf{Metoda} & \textbf{Endpoint} & \textbf{Opis} \\
\hline
\endhead
\multicolumn{3}{|c|}{\textbf{Autentykacja}} \\
\hline
POST & /api/auth/register & Rejestracja użytkownika \\
POST & /api/auth/login & Logowanie \\
POST & /api/auth/refresh & Odświeżenie tokenu \\
POST & /api/auth/forgot-password & Żądanie resetu hasła \\
POST & /api/auth/reset-password & Reset hasła \\
\hline
\multicolumn{3}{|c|}{\textbf{Projekty}} \\
\hline
GET & /api/projects & Lista projektów \\
POST & /api/projects & Tworzenie projektu \\
GET & /api/projects/\{id\} & Szczegóły projektu \\
PUT & /api/projects/\{id\} & Aktualizacja projektu \\
DELETE & /api/projects/\{id\} & Usunięcie projektu \\
\hline
\multicolumn{3}{|c|}{\textbf{Zadania}} \\
\hline
GET & /api/tasks & Lista zadań \\
POST & /api/tasks & Tworzenie zadania \\
GET & /api/tasks/\{id\} & Szczegóły zadania \\
PUT & /api/tasks/\{id\} & Aktualizacja zadania \\
DELETE & /api/tasks/\{id\} & Usunięcie zadania \\
GET & /api/tasks/project/\{projectId\} & Zadania projektu \\
\hline
\multicolumn{3}{|c|}{\textbf{Powiadomienia}} \\
\hline
GET & /api/notifications & Lista powiadomień \\
GET & /api/notifications/unread & Nieprzeczytane \\
PUT & /api/notifications/\{id\}/read & Oznacz jako przeczytane \\
PUT & /api/notifications/read-all & Oznacz wszystkie \\
DELETE & /api/notifications/\{id\} & Usuń powiadomienie \\
\hline
\multicolumn{3}{|c|}{\textbf{Audyt (Admin)}} \\
\hline
GET & /api/audit & Wszystkie logi \\
GET & /api/audit/user/\{userId\} & Logi użytkownika \\
GET & /api/audit/entity/\{entity\}/\{id\} & Logi encji \\
\hline
\end{longtable}

\subsection{Przykładowe formularze (DTO)}

\begin{lstlisting}[style=javastyle, caption={RegisterRequest - formularz rejestracji}]
@Data
@Builder
public class RegisterRequest {
    @NotBlank(message = "Imie jest wymagane")
    @Size(min = 1, max = 50)
    private String firstName;
    
    @NotBlank(message = "Nazwisko jest wymagane")
    @Size(min = 1, max = 50)
    private String lastName;
    
    @NotBlank(message = "Email jest wymagany")
    @Email(message = "Nieprawidlowy format email")
    private String email;
    
    @NotBlank(message = "Haslo jest wymagane")
    @Size(min = 6, message = "Haslo min. 6 znakow")
    private String password;
}
\end{lstlisting}

\begin{lstlisting}[style=javastyle, caption={ProjectRequest - formularz projektu}]
@Data
@Builder
public class ProjectRequest {
    @NotBlank(message = "Nazwa projektu jest wymagana")
    @Size(min = 1, max = 255)
    private String name;
    
    @Size(max = 1000)
    private String description;
    
    private LocalDate startDate;
    private LocalDate endDate;
    private ProjectStatus status;
}
\end{lstlisting}

% ============================================================================
% ROZDZIAŁ 3: NORMALIZACJA I WIĘZY
% ============================================================================
\chapter{Normalizacja i więzy tabel}

\section{Analiza normalizacji}

\subsection{Pierwsza Postać Normalna (1NF)}

Wszystkie tabele w systemie spełniają wymagania 1NF:

\begin{itemize}
    \item Każda kolumna zawiera wartości atomowe (niepodzielne)
    \item Każda tabela posiada klucz główny (PRIMARY KEY)
    \item Brak powtarzających się grup kolumn
    \item Każda komórka zawiera pojedynczą wartość
\end{itemize}

\textbf{Przykład -- tabela USERS:}
\begin{itemize}
    \item Imię i nazwisko rozdzielone na osobne kolumny (first\_name, last\_name)
    \item Email jako pojedyncza wartość
    \item Brak list ani tablic w kolumnach
\end{itemize}

\subsection{Druga Postać Normalna (2NF)}

Wszystkie tabele spełniają 2NF:

\begin{itemize}
    \item Spełniają 1NF
    \item Wszystkie atrybuty niekluczowe są w pełni funkcyjnie zależne od klucza głównego
    \item Brak częściowych zależności funkcyjnych
\end{itemize}

\textbf{Przykład -- tabela PROJECT\_USER:}
\begin{itemize}
    \item Klucz złożony: (project\_id, user\_id)
    \item Atrybut \texttt{role} zależy od całego klucza złożonego
    \item Nie można określić roli znając tylko projekt lub tylko użytkownika
\end{itemize}

\subsection{Trzecia Postać Normalna (3NF)}

Wszystkie tabele spełniają 3NF:

\begin{itemize}
    \item Spełniają 2NF
    \item Brak zależności przechodnich
    \item Każdy atrybut niekluczowy zależy bezpośrednio od klucza głównego
\end{itemize}

\textbf{Przykład -- tabela TASKS:}
\begin{itemize}
    \item Wszystkie atrybuty (name, description, status, due\_date) zależą bezpośrednio od id
    \item project\_id i assigned\_to to klucze obce, nie powodują zależności przechodnich
    \item Dane projektu nie są duplikowane w tabeli zadań
\end{itemize}

\subsection{Postać Normalna Boyce'a-Codda (BCNF)}

Tabela PROJECT\_USER jest w BCNF:
\begin{itemize}
    \item Każdy wyznacznik jest kluczem kandydującym
    \item Klucz złożony (project\_id, user\_id) jest jedynym wyznacznikiem
\end{itemize}

\section{Więzy integralności}

\subsection{Więzy klucza głównego (PRIMARY KEY)}

\begin{lstlisting}[style=sqlstyle, caption={Definicje kluczy głównych}]
-- Klucz prosty
ALTER TABLE users ADD CONSTRAINT pk_users PRIMARY KEY (id);
ALTER TABLE projects ADD CONSTRAINT pk_projects PRIMARY KEY (id);
ALTER TABLE tasks ADD CONSTRAINT pk_tasks PRIMARY KEY (id);

-- Klucz zlozony
ALTER TABLE project_user 
ADD CONSTRAINT pk_project_user PRIMARY KEY (project_id, user_id);
\end{lstlisting}

\subsection{Więzy klucza obcego (FOREIGN KEY)}

\begin{lstlisting}[style=sqlstyle, caption={Definicje kluczy obcych}]
-- Projects -> Users (tworca)
ALTER TABLE projects 
ADD CONSTRAINT fk_project_creator 
FOREIGN KEY (created_by) REFERENCES users(id)
ON DELETE RESTRICT;

-- Tasks -> Projects
ALTER TABLE tasks 
ADD CONSTRAINT fk_task_project 
FOREIGN KEY (project_id) REFERENCES projects(id)
ON DELETE CASCADE;

-- Tasks -> Users (przypisanie)
ALTER TABLE tasks 
ADD CONSTRAINT fk_task_assignee 
FOREIGN KEY (assigned_to) REFERENCES users(id)
ON DELETE SET NULL;

-- Project_User -> Projects
ALTER TABLE project_user 
ADD CONSTRAINT fk_pu_project 
FOREIGN KEY (project_id) REFERENCES projects(id)
ON DELETE CASCADE;

-- Project_User -> Users
ALTER TABLE project_user 
ADD CONSTRAINT fk_pu_user 
FOREIGN KEY (user_id) REFERENCES users(id)
ON DELETE CASCADE;
\end{lstlisting}

\subsection{Więzy unikalności (UNIQUE)}

\begin{lstlisting}[style=sqlstyle, caption={Więzy UNIQUE}]
-- Unikalny email uzytkownika
ALTER TABLE users ADD CONSTRAINT uq_user_email UNIQUE (email);

-- Unikalna nazwa przechowywanego pliku
ALTER TABLE project_files 
ADD CONSTRAINT uq_stored_filename UNIQUE (stored_file_name);

-- Unikalny token odswiezania
ALTER TABLE refresh_tokens ADD CONSTRAINT uq_token UNIQUE (token);
\end{lstlisting}

\subsection{Więzy CHECK}

\begin{lstlisting}[style=sqlstyle, caption={Więzy CHECK}]
-- Walidacja statusu projektu
ALTER TABLE projects ADD CONSTRAINT chk_project_status 
CHECK (status IN ('NOT_STARTED', 'IN_PROGRESS', 
                  'ON_HOLD', 'COMPLETED', 'CANCELLED'));

-- Walidacja statusu zadania
ALTER TABLE tasks ADD CONSTRAINT chk_task_status 
CHECK (status IN ('TODO', 'IN_PROGRESS', 'REVIEW', 'DONE'));

-- Walidacja roli uzytkownika
ALTER TABLE users ADD CONSTRAINT chk_user_role 
CHECK (role IN ('USER', 'ADMIN'));

-- Walidacja roli w projekcie
ALTER TABLE project_user ADD CONSTRAINT chk_project_role 
CHECK (role IN ('OWNER', 'MEMBER', 'VIEWER'));

-- Walidacja rozmiaru pliku
ALTER TABLE project_files ADD CONSTRAINT chk_file_size 
CHECK (file_size > 0);

-- Walidacja dat projektu
ALTER TABLE projects ADD CONSTRAINT chk_project_dates 
CHECK (end_date IS NULL OR start_date IS NULL 
       OR end_date >= start_date);
\end{lstlisting}

\subsection{Więzy NOT NULL}

\begin{lstlisting}[style=sqlstyle, caption={Więzy NOT NULL}]
-- Tabela users
ALTER TABLE users ALTER COLUMN first_name SET NOT NULL;
ALTER TABLE users ALTER COLUMN last_name SET NOT NULL;
ALTER TABLE users ALTER COLUMN email SET NOT NULL;
ALTER TABLE users ALTER COLUMN password SET NOT NULL;
ALTER TABLE users ALTER COLUMN role SET NOT NULL;

-- Tabela projects
ALTER TABLE projects ALTER COLUMN name SET NOT NULL;
ALTER TABLE projects ALTER COLUMN status SET NOT NULL;
ALTER TABLE projects ALTER COLUMN created_by SET NOT NULL;

-- Tabela tasks
ALTER TABLE tasks ALTER COLUMN name SET NOT NULL;
ALTER TABLE tasks ALTER COLUMN status SET NOT NULL;
ALTER TABLE tasks ALTER COLUMN project_id SET NOT NULL;
\end{lstlisting}

\section{Akcje referencyjne}

\begin{table}[H]
\centering
\caption{Akcje referencyjne dla kluczy obcych}
\begin{tabular}{|l|l|l|l|}
\hline
\textbf{Tabela} & \textbf{FK} & \textbf{ON DELETE} & \textbf{Uzasadnienie} \\
\hline
projects & created\_by & RESTRICT & Nie usuwaj użytkownika z projektami \\
tasks & project\_id & CASCADE & Usuń zadania z projektem \\
tasks & assigned\_to & SET NULL & Zachowaj zadanie, usuń przypisanie \\
project\_user & project\_id & CASCADE & Usuń powiązania z projektem \\
project\_user & user\_id & CASCADE & Usuń powiązania użytkownika \\
comments & project\_id & CASCADE & Usuń komentarze z projektem \\
comments & user\_id & CASCADE & Usuń komentarze użytkownika \\
files & project\_id & CASCADE & Usuń pliki z projektem \\
notifications & user\_id & CASCADE & Usuń powiadomienia użytkownika \\
audit\_logs & user\_id & CASCADE & Usuń logi użytkownika \\
\hline
\end{tabular}
\end{table}

% ============================================================================
% ROZDZIAŁ 4: KONTROLA POPRAWNOŚCI DANYCH
% ============================================================================
\chapter{Kontrola poprawności danych}

\section{Walidacja na poziomie encji (Bean Validation)}

System wykorzystuje Jakarta Bean Validation (JSR-380) do walidacji danych na poziomie modelu.

\begin{lstlisting}[style=javastyle, caption={Walidacja w encji User}]
@Entity
@Table(name = "users")
public class User {

    @NotBlank(message = "Imie jest wymagane")
    @Size(min = 1, max = 50, 
          message = "Imie musi miec od 1 do 50 znakow")
    @Column(nullable = false, length = 50)
    private String firstName;

    @NotBlank(message = "Nazwisko jest wymagane")
    @Size(min = 1, max = 50, 
          message = "Nazwisko musi miec od 1 do 50 znakow")
    @Column(nullable = false, length = 50)
    private String lastName;

    @NotBlank(message = "Email jest wymagany")
    @Email(message = "Nieprawidlowy format email")
    @Column(nullable = false, unique = true, length = 100)
    private String email;
}
\end{lstlisting}

\begin{lstlisting}[style=javastyle, caption={Walidacja w encji Project}]
@Entity
@Table(name = "projects")
public class Project {

    @NotBlank(message = "Nazwa projektu jest wymagana")
    @Size(min = 1, max = 255, 
          message = "Nazwa projektu musi miec od 1 do 255 znakow")
    @Column(nullable = false, length = 255)
    private String name;

    @Size(max = 1000, 
          message = "Opis nie moze przekraczac 1000 znakow")
    @Column(length = 1000)
    private String description;

    // Walidacja dat w metodzie @PrePersist/@PreUpdate
    @PrePersist
    @PreUpdate
    private void validateDates() {
        if (startDate != null && endDate != null 
            && endDate.isBefore(startDate)) {
            throw new IllegalArgumentException(
                "Data zakonczenia nie moze byc 
                 wczesniejsza niz data rozpoczecia");
        }
    }
}
\end{lstlisting}

\section{Walidacja w kontrolerach}

\begin{lstlisting}[style=javastyle, caption={Walidacja w kontrolerze z @Valid}]
@RestController
@RequestMapping("/api/projects")
public class ProjectController {

    @PostMapping
    public ResponseEntity<ProjectResponse> createProject(
            @Valid @RequestBody ProjectRequest request,
            Authentication authentication) {
        // @Valid wymusza walidacje przed wykonaniem metody
        ProjectResponse response = projectService
            .createProject(request, authentication.getName());
        return ResponseEntity.status(HttpStatus.CREATED)
            .body(response);
    }
}
\end{lstlisting}

\section{Obsługa wyjątków}

\begin{lstlisting}[style=javastyle, caption={GlobalExceptionHandler}]
@RestControllerAdvice
public class GlobalExceptionHandler {

    // Obsluga bledow walidacji
    @ExceptionHandler(MethodArgumentNotValidException.class)
    public ResponseEntity<Map<String, String>> 
            handleValidationExceptions(
                MethodArgumentNotValidException ex) {
        Map<String, String> errors = new HashMap<>();
        ex.getBindingResult().getAllErrors().forEach(error -> {
            String fieldName = ((FieldError) error).getField();
            String errorMessage = error.getDefaultMessage();
            errors.put(fieldName, errorMessage);
        });
        return ResponseEntity.badRequest().body(errors);
    }

    // Obsluga ResourceNotFoundException
    @ExceptionHandler(ResourceNotFoundException.class)
    public ResponseEntity<ErrorResponse> 
            handleResourceNotFound(ResourceNotFoundException ex) {
        ErrorResponse error = new ErrorResponse(
            HttpStatus.NOT_FOUND.value(),
            ex.getMessage(),
            LocalDateTime.now()
        );
        return ResponseEntity.status(HttpStatus.NOT_FOUND)
            .body(error);
    }

    // Obsluga UnauthorizedOperationException
    @ExceptionHandler(UnauthorizedOperationException.class)
    public ResponseEntity<ErrorResponse> 
            handleUnauthorized(UnauthorizedOperationException ex) {
        ErrorResponse error = new ErrorResponse(
            HttpStatus.FORBIDDEN.value(),
            ex.getMessage(),
            LocalDateTime.now()
        );
        return ResponseEntity.status(HttpStatus.FORBIDDEN)
            .body(error);
    }
}
\end{lstlisting}

\section{Wartości domyślne}

\begin{lstlisting}[style=javastyle, caption={Wartości domyślne w encjach}]
@Entity
public class Project {
    // Domyslny status projektu
    @Enumerated(EnumType.STRING)
    @Builder.Default
    private ProjectStatus status = ProjectStatus.NOT_STARTED;
}

@Entity
public class Notification {
    // Domyslna wartosc "przeczytane"
    @Column(nullable = false)
    @Builder.Default
    private boolean read = false;
}

@Entity
public class User {
    // Automatyczna data utworzenia
    @CreationTimestamp
    @Column(updatable = false)
    private LocalDateTime createdAt;
}
\end{lstlisting}

\section{Walidacja plików}

\begin{lstlisting}[style=javastyle, caption={Walidacja uploadowanych plików}]
@Service
public class FileStorageService {

    private static final List<String> ALLOWED_TYPES = Arrays.asList(
        "image/jpeg", "image/png", "image/gif",
        "application/pdf", "application/msword",
        "application/vnd.openxmlformats-officedocument"
    );

    @Value("${file.max-file-size-mb}")
    private int maxFileSizeMb;

    public String storeFile(MultipartFile file, String subDir) {
        // Walidacja rozmiaru
        if (file.getSize() > maxFileSizeMb * 1024 * 1024) {
            throw new FileStorageException(
                "Plik przekracza maksymalny rozmiar " 
                + maxFileSizeMb + "MB");
        }

        // Walidacja typu
        String contentType = file.getContentType();
        if (contentType == null || 
            !ALLOWED_TYPES.contains(contentType)) {
            throw new FileStorageException(
                "Niedozwolony typ pliku: " + contentType);
        }

        // Walidacja nazwy pliku
        String originalFilename = StringUtils
            .cleanPath(file.getOriginalFilename());
        if (originalFilename.contains("..")) {
            throw new FileStorageException(
                "Nazwa pliku zawiera nieprawidlowa sciezke");
        }

        // Generowanie unikalnej nazwy
        String storedFileName = UUID.randomUUID().toString() 
            + "_" + originalFilename;
        
        // Zapisanie pliku...
        return storedFileName;
    }
}
\end{lstlisting}

\section{Maski wprowadzania (wzorce regex)}

\begin{lstlisting}[style=javastyle, caption={Wzorce walidacji}]
public class ValidationPatterns {
    
    // Wzorzec email
    public static final String EMAIL_PATTERN = 
        "^[A-Za-z0-9._%+-]+@[A-Za-z0-9.-]+\\.[A-Za-z]{2,}$";
    
    // Wzorzec hasla (min 8 znakow, litera, cyfra)
    public static final String PASSWORD_PATTERN = 
        "^(?=.*[A-Za-z])(?=.*\\d)[A-Za-z\\d@$!%*#?&]{8,}$";
    
    // Wzorzec nazwy uzytkownika
    public static final String USERNAME_PATTERN = 
        "^[a-zA-Z][a-zA-Z0-9_]{2,29}$";
}

// Uzycie w DTO
public class RegisterRequest {
    @Pattern(regexp = ValidationPatterns.EMAIL_PATTERN,
             message = "Nieprawidlowy format email")
    private String email;
    
    @Pattern(regexp = ValidationPatterns.PASSWORD_PATTERN,
             message = "Haslo musi zawierac min. 8 znakow, "
                     + "litere i cyfre")
    private String password;
}
\end{lstlisting}

% ============================================================================
% ROZDZIAŁ 5: AUTOMATYCZNE TWORZENIE TABEL
% ============================================================================
\chapter{Automatyczne tworzenie tabel}

\section{Konfiguracja Hibernate DDL}

System wykorzystuje Hibernate do automatycznego tworzenia i aktualizacji schematu bazy danych.

\begin{lstlisting}[style=propstyle, caption={application.properties - konfiguracja DDL}]
# Konfiguracja JPA/Hibernate
spring.jpa.hibernate.ddl-auto=create-drop
spring.jpa.show-sql=true
spring.jpa.properties.hibernate.format_sql=true
spring.jpa.properties.hibernate.dialect=
    org.hibernate.dialect.PostgreSQLDialect

# Dla srodowiska produkcyjnego:
# spring.jpa.hibernate.ddl-auto=validate
\end{lstlisting}

\subsection{Dostępne tryby ddl-auto}

\begin{table}[H]
\centering
\caption{Tryby ddl-auto w Hibernate}
\begin{tabular}{|l|p{10cm}|}
\hline
\textbf{Tryb} & \textbf{Opis} \\
\hline
\texttt{none} & Brak automatycznego zarządzania schematem \\
\texttt{validate} & Waliduje schemat, nie wprowadza zmian \\
\texttt{update} & Aktualizuje schemat (dodaje brakujące elementy) \\
\texttt{create} & Tworzy schemat przy starcie (usuwa istniejący) \\
\texttt{create-drop} & Tworzy przy starcie, usuwa przy zamknięciu \\
\hline
\end{tabular}
\end{table}

\section{Definicje encji JPA}

\begin{lstlisting}[style=javastyle, caption={Encja Project z adnotacjami JPA}]
@Entity
@Table(name = "projects", indexes = {
    @Index(name = "idx_project_created_by", 
           columnList = "created_by"),
    @Index(name = "idx_project_status", 
           columnList = "status"),
    @Index(name = "idx_project_created_at", 
           columnList = "createdAt"),
    @Index(name = "idx_project_name", 
           columnList = "name")
})
@Getter
@Setter
@NoArgsConstructor
@AllArgsConstructor
@Builder
public class Project {

    @Id
    @GeneratedValue(strategy = GenerationType.IDENTITY)
    private Long id;

    @Column(nullable = false, length = 255)
    private String name;

    @Column(length = 1000)
    private String description;

    @Enumerated(EnumType.STRING)
    @Column(nullable = false, length = 20)
    @Builder.Default
    private ProjectStatus status = ProjectStatus.NOT_STARTED;

    @ManyToOne(fetch = FetchType.LAZY)
    @JoinColumn(name = "created_by", nullable = false)
    private User createdBy;

    // Relacje z kaskadowym usuwaniem
    @OneToMany(mappedBy = "project", 
               cascade = CascadeType.ALL, 
               orphanRemoval = true)
    @Builder.Default
    private List<Task> tasks = new ArrayList<>();

    @OneToMany(mappedBy = "project", 
               cascade = CascadeType.ALL, 
               orphanRemoval = true)
    @Builder.Default
    private List<ProjectComment> comments = new ArrayList<>();
}
\end{lstlisting}

\section{Inicjalizacja danych testowych}

\begin{lstlisting}[style=javastyle, caption={DatabaseInitializer - inicjalizacja danych}]
@Component
@RequiredArgsConstructor
public class DatabaseInitializer implements CommandLineRunner {

    private final UserRepository userRepository;
    private final PasswordEncoder passwordEncoder;
    private final ProjectRepository projectRepository;
    private final TaskRepository taskRepository;

    @Override
    public void run(String... args) {
        // Sprawdz czy baza jest pusta
        if (userRepository.count() == 0) {
            initUsers();
        }

        if (projectRepository.count() == 0) {
            initProjects();
        }

        if (taskRepository.count() == 0) {
            initTasks();
        }
    }

    private void initUsers() {
        List<User> users = Arrays.asList(
            User.builder()
                .firstName("Admin")
                .lastName("Systemowy")
                .email("admin@example.com")
                .password(passwordEncoder.encode("admin123"))
                .role(Role.ADMIN)
                .build(),
            User.builder()
                .firstName("Jan")
                .lastName("Kowalski")
                .email("jan@example.com")
                .password(passwordEncoder.encode("user123"))
                .role(Role.USER)
                .build()
        );
        userRepository.saveAll(users);
    }
}
\end{lstlisting}

% ============================================================================
% ROZDZIAŁ 6: REJESTRACJA DATY I CZASU
% ============================================================================
\chapter{Rejestracja daty i czasu operacji}

\section{Automatyczne znaczniki czasowe}

System wykorzystuje adnotacje Hibernate do automatycznego rejestrowania znaczników czasowych.

\begin{lstlisting}[style=javastyle, caption={Automatyczne timestampy w encjach}]
@Entity
public class User {
    // Automatyczna data utworzenia - niemodyfikowalna
    @CreationTimestamp
    @Column(updatable = false)
    private LocalDateTime createdAt;
}

@Entity
public class AuditLog {
    // Automatyczny timestamp operacji
    @CreationTimestamp
    @Column(updatable = false, nullable = false)
    private LocalDateTime timestamp;
}

@Entity
public class ProjectComment {
    // Data utworzenia
    @CreationTimestamp
    @Column(updatable = false)
    private LocalDateTime createdAt;

    // Data modyfikacji - reczna aktualizacja
    @Column
    private LocalDateTime updatedAt;

    @PreUpdate
    protected void onUpdate() {
        updatedAt = LocalDateTime.now();
    }
}
\end{lstlisting}

\section{System audytu}

\begin{lstlisting}[style=javastyle, caption={AuditLogService - logowanie operacji}]
@Service
@RequiredArgsConstructor
public class AuditLogService {

    private final AuditLogRepository auditLogRepository;
    private final UserRepository userRepository;

    // Logowanie akcji asynchronicznie
    @Async
    @Transactional
    public void logActionAsync(String userEmail, String action, 
                               String entity, Long entityId) {
        logAction(userEmail, action, entity, entityId);
    }

    // Logowanie akcji synchronicznie
    @Transactional
    public AuditLog logAction(String userEmail, String action, 
                              String entity, Long entityId) {
        User user = userRepository.findByEmail(userEmail)
            .orElseThrow(() -> 
                new ResourceNotFoundException("User", "email", userEmail));
        
        AuditLog auditLog = AuditLog.builder()
            .user(user)
            .action(action)
            .entity(entity)
            .entityId(entityId)
            .ipAddress(SecurityUtils.getClientIpAddress())
            .build();
        
        return auditLogRepository.save(auditLog);
    }

    // Stale dla typow akcji
    public static final String ACTION_CREATE = "CREATE";
    public static final String ACTION_UPDATE = "UPDATE";
    public static final String ACTION_DELETE = "DELETE";
    public static final String ACTION_VIEW = "VIEW";
}
\end{lstlisting}

\section{Integracja audytu z serwisami}

\begin{lstlisting}[style=javastyle, caption={Użycie audytu w ProjectService}]
@Service
@RequiredArgsConstructor
public class ProjectService {

    private final AuditLogService auditLogService;

    @Transactional
    public ProjectResponse createProject(
            ProjectRequest request, String userEmail) {
        // ... tworzenie projektu ...
        
        Project savedProject = projectRepository.save(project);
        
        // Logowanie audytu
        auditLogService.logActionAsync(
            userEmail, 
            AuditLogService.ACTION_CREATE, 
            AuditLogService.ENTITY_PROJECT, 
            savedProject.getId()
        );
        
        return convertToResponse(savedProject);
    }

    @Transactional
    public void deleteProject(Long projectId, String userEmail) {
        // ... walidacja ...
        
        // Logowanie PRZED usunieciem (synchronicznie)
        auditLogService.logAction(
            userEmail, 
            AuditLogService.ACTION_DELETE, 
            AuditLogService.ENTITY_PROJECT, 
            projectId
        );
        
        projectRepository.delete(project);
    }
}
\end{lstlisting}

\section{Możliwość ręcznej korekty}

\begin{lstlisting}[style=javastyle, caption={Ręczna korekta timestamp (Admin)}]
@RestController
@RequestMapping("/api/admin/audit")
@PreAuthorize("hasRole('ADMIN')")
public class AdminAuditController {

    @PutMapping("/{logId}/timestamp")
    public ResponseEntity<AuditLogResponse> correctTimestamp(
            @PathVariable Long logId,
            @RequestBody TimestampCorrectionRequest request) {
        
        AuditLog log = auditLogRepository.findById(logId)
            .orElseThrow(() -> 
                new ResourceNotFoundException("AuditLog", "id", logId));
        
        // Reczna korekta timestamp (wymaga uprawnien ADMIN)
        // Pole timestamp w AuditLog musi byc bez "updatable = false"
        log.setTimestamp(request.getNewTimestamp());
        
        // Loguj korekte
        auditLogService.logAction(
            getCurrentUserEmail(),
            "TIMESTAMP_CORRECTION",
            "AuditLog",
            logId
        );
        
        return ResponseEntity.ok(convertToResponse(
            auditLogRepository.save(log)));
    }
}
\end{lstlisting}

% ============================================================================
% ROZDZIAŁ 7: INSTRUKCJA OBSŁUGI
% ============================================================================
\chapter{Instrukcja obsługi}

\section{Wymagania systemowe}

\begin{table}[H]
\centering
\caption{Wymagania systemowe}
\begin{tabular}{|l|l|}
\hline
\textbf{Komponent} & \textbf{Wersja minimalna} \\
\hline
Java JDK & 23 \\
PostgreSQL & 15 \\
Gradle & 8.0 \\
Docker (opcjonalnie) & 20.0 \\
\hline
\end{tabular}
\end{table}

\section{Instalacja}

\subsection{Klonowanie repozytorium}

\begin{lstlisting}[language=bash, caption={Klonowanie projektu}]
git clone https://github.com/puchacz51/zw-backend.git
cd zw-backend
\end{lstlisting}

\subsection{Konfiguracja bazy danych}

\begin{lstlisting}[style=propstyle, caption={application.properties}]
# Konfiguracja polaczenia z baza danych
spring.datasource.url=jdbc:postgresql://localhost:5432/mydatabase
spring.datasource.username=user
spring.datasource.password=secret
spring.datasource.driver-class-name=org.postgresql.Driver

# Konfiguracja JWT
jwt.secret=twojBardzoTajnyKlucz256bitow
jwt.accessTokenExpirationMs=3600000
jwt.refreshTokenExpirationMs=86400000

# Konfiguracja plikow
file.base-upload-dir=./uploads
file.max-file-size-mb=50
\end{lstlisting}

\subsection{Uruchomienie z Docker Compose}

\begin{lstlisting}[language=bash, caption={Uruchomienie z Docker}]
# Uruchomienie bazy danych i aplikacji
docker-compose up -d

# Sprawdzenie statusu
docker-compose ps

# Logi aplikacji
docker-compose logs -f app
\end{lstlisting}

\subsection{Uruchomienie lokalne}

\begin{lstlisting}[language=bash, caption={Uruchomienie lokalne}]
# Kompilacja projektu
./gradlew build

# Uruchomienie aplikacji
./gradlew bootRun

# Lub uruchomienie JAR
java -jar build/libs/zw-backend-0.0.1-SNAPSHOT.jar
\end{lstlisting}

\section{Dokumentacja API}

Po uruchomieniu aplikacji, dokumentacja API dostępna jest pod adresami:

\begin{itemize}
    \item \textbf{Swagger UI:} \url{http://localhost:8080/swagger-ui.html}
    \item \textbf{OpenAPI JSON:} \url{http://localhost:8080/v3/api-docs}
\end{itemize}

\section{Przykłady użycia API}

\subsection{Rejestracja użytkownika}

\begin{lstlisting}[language=bash, caption={curl - rejestracja}]
curl -X POST http://localhost:8080/api/auth/register \
  -H "Content-Type: application/json" \
  -d '{
    "firstName": "Jan",
    "lastName": "Kowalski",
    "email": "jan@example.com",
    "password": "haslo123"
  }'
\end{lstlisting}

\subsection{Logowanie}

\begin{lstlisting}[language=bash, caption={curl - logowanie}]
curl -X POST http://localhost:8080/api/auth/login \
  -H "Content-Type: application/json" \
  -d '{
    "email": "jan@example.com",
    "password": "haslo123"
  }'
\end{lstlisting}

\subsection{Tworzenie projektu}

\begin{lstlisting}[language=bash, caption={curl - tworzenie projektu}]
curl -X POST http://localhost:8080/api/projects \
  -H "Content-Type: application/json" \
  -H "Authorization: Bearer <JWT_TOKEN>" \
  -d '{
    "name": "Nowy Projekt",
    "description": "Opis projektu",
    "startDate": "2026-01-28",
    "endDate": "2026-06-30"
  }'
\end{lstlisting}

\section{Uruchomienie testów}

\begin{lstlisting}[language=bash, caption={Uruchomienie testów}]
# Wszystkie testy
./gradlew test

# Testy jednostkowe
./gradlew test --tests "*Test"

# Testy integracyjne
./gradlew test --tests "*IntegrationTest"

# Raport z testow
open build/reports/tests/test/index.html
\end{lstlisting}

% ============================================================================
% ROZDZIAŁ 8: KOD ŹRÓDŁOWY
% ============================================================================
\chapter{Kod źródłowy -- wybrane fragmenty}

\section{Struktura projektu}

\begin{lstlisting}[language=bash, caption={Struktura katalogów}]
src/
|-- main/
|   |-- java/pl/pbs/zwbackend/
|   |   |-- ZwBackendApplication.java
|   |   |-- config/
|   |   |   |-- SecurityConfig.java
|   |   |   |-- OpenApiConfig.java
|   |   |   |-- AsyncConfig.java
|   |   |-- controller/
|   |   |   |-- AuthController.java
|   |   |   |-- ProjectController.java
|   |   |   |-- TaskController.java
|   |   |   |-- NotificationController.java
|   |   |   |-- AuditLogController.java
|   |   |-- dto/
|   |   |-- exception/
|   |   |-- model/
|   |   |   |-- User.java
|   |   |   |-- Project.java
|   |   |   |-- Task.java
|   |   |   |-- enums/
|   |   |-- repository/
|   |   |-- security/
|   |   |-- service/
|   |   |-- util/
|   |-- resources/
|       |-- application.properties
|-- test/
    |-- java/pl/pbs/zwbackend/
        |-- controller/
        |-- integration/
        |-- repository/
        |-- service/
\end{lstlisting}

\section{Główna klasa aplikacji}

\begin{lstlisting}[style=javastyle, caption={ZwBackendApplication.java}]
package pl.pbs.zwbackend;

import org.springframework.boot.SpringApplication;
import org.springframework.boot.autoconfigure.SpringBootApplication;
import org.springframework.scheduling.annotation.EnableAsync;

@SpringBootApplication
@EnableAsync
public class ZwBackendApplication {

    public static void main(String[] args) {
        SpringApplication.run(ZwBackendApplication.class, args);
    }
}
\end{lstlisting}

\section{Konfiguracja bezpieczeństwa}

\begin{lstlisting}[style=javastyle, caption={SecurityConfig.java}]
@Configuration
@EnableWebSecurity
@EnableMethodSecurity
@RequiredArgsConstructor
public class SecurityConfig {

    private final JwtAuthenticationFilter jwtAuthFilter;
    private final CustomUserDetailsService userDetailsService;

    @Bean
    public SecurityFilterChain securityFilterChain(
            HttpSecurity http) throws Exception {
        http
            .csrf(csrf -> csrf.disable())
            .cors(cors -> cors.configurationSource(
                corsConfigurationSource()))
            .sessionManagement(session -> session
                .sessionCreationPolicy(
                    SessionCreationPolicy.STATELESS))
            .authorizeHttpRequests(auth -> auth
                .requestMatchers("/api/auth/**").permitAll()
                .requestMatchers("/swagger-ui/**").permitAll()
                .requestMatchers("/v3/api-docs/**").permitAll()
                .requestMatchers("/api/admin/**")
                    .hasRole("ADMIN")
                .anyRequest().authenticated()
            )
            .authenticationProvider(authenticationProvider())
            .addFilterBefore(jwtAuthFilter, 
                UsernamePasswordAuthenticationFilter.class);
        
        return http.build();
    }

    @Bean
    public PasswordEncoder passwordEncoder() {
        return new BCryptPasswordEncoder();
    }
}
\end{lstlisting}

\section{Kontroler autentykacji}

\begin{lstlisting}[style=javastyle, caption={AuthController.java}]
@RestController
@RequestMapping("/api/auth")
@RequiredArgsConstructor
@Tag(name = "Authentication", 
     description = "API do autentykacji")
public class AuthController {

    private final AuthenticationManager authManager;
    private final JwtTokenProvider jwtProvider;
    private final UserService userService;

    @PostMapping("/register")
    @Operation(summary = "Rejestracja nowego uzytkownika")
    public ResponseEntity<UserSummaryResponse> register(
            @Valid @RequestBody RegisterRequest request) {
        User user = userService.createUser(request);
        return ResponseEntity.status(HttpStatus.CREATED)
            .body(userService.convertToUserSummaryResponse(user));
    }

    @PostMapping("/login")
    @Operation(summary = "Logowanie uzytkownika")
    public ResponseEntity<TokenResponse> login(
            @Valid @RequestBody LoginRequest request) {
        Authentication auth = authManager.authenticate(
            new UsernamePasswordAuthenticationToken(
                request.getEmail(), 
                request.getPassword()
            )
        );
        
        String accessToken = jwtProvider
            .generateAccessToken(auth);
        String refreshToken = jwtProvider
            .generateRefreshToken(auth);
        
        return ResponseEntity.ok(TokenResponse.builder()
            .accessToken(accessToken)
            .refreshToken(refreshToken)
            .tokenType("Bearer")
            .build());
    }
}
\end{lstlisting}

% ============================================================================
% ROZDZIAŁ 9: PREZENTACJA PROJEKTU
% ============================================================================
\chapter{Prezentacja projektu}

\section{Scenariusze testowe}

\subsection{Scenariusz 1: Rejestracja i logowanie}

\begin{enumerate}
    \item Użytkownik otwiera aplikację
    \item Wypełnia formularz rejestracji (imię, nazwisko, email, hasło)
    \item System waliduje dane i tworzy konto
    \item Użytkownik loguje się podając email i hasło
    \item System zwraca token JWT
    \item Użytkownik używa tokenu do autoryzacji żądań
\end{enumerate}

\subsection{Scenariusz 2: Zarządzanie projektem}

\begin{enumerate}
    \item Zalogowany użytkownik tworzy nowy projekt
    \item Dodaje opis, daty rozpoczęcia i zakończenia
    \item Przypisuje innych użytkowników do projektu
    \item Tworzy zadania w ramach projektu
    \item Przypisuje zadania do członków zespołu
    \item Śledzi postęp przez zmianę statusów
\end{enumerate}

\subsection{Scenariusz 3: Współpraca zespołowa}

\begin{enumerate}
    \item Członek zespołu otrzymuje powiadomienie o przypisaniu
    \item Przegląda szczegóły zadania
    \item Dodaje komentarz z pytaniem
    \item Przesyła plik z wynikami pracy
    \item Zmienia status zadania na "W trakcie"
    \item Po zakończeniu zmienia status na "Gotowe"
\end{enumerate}

\section{Demonstracja funkcjonalności}

\subsection{Ekran logowania}

\begin{figure}[H]
\centering
\fbox{\parbox{0.8\textwidth}{
\centering
\vspace{2cm}
\textbf{[Miejsce na zrzut ekranu]}\\
Formularz logowania z polami email i hasło
\vspace{2cm}
}}
\caption{Ekran logowania (Swagger UI)}
\end{figure}

\subsection{Lista projektów}

\begin{figure}[H]
\centering
\fbox{\parbox{0.8\textwidth}{
\centering
\vspace{2cm}
\textbf{[Miejsce na zrzut ekranu]}\\
Lista projektów z statusami i członkami zespołu
\vspace{2cm}
}}
\caption{Lista projektów}
\end{figure}

\section{Wykorzystane narzędzia}

\begin{table}[H]
\centering
\caption{Narzędzia wykorzystane w projekcie}
\begin{tabular}{|l|l|}
\hline
\textbf{Kategoria} & \textbf{Narzędzie} \\
\hline
IDE & IntelliJ IDEA / VS Code \\
Kontrola wersji & Git, GitHub \\
Baza danych & PostgreSQL, pgAdmin \\
Testowanie API & Postman, Swagger UI \\
Konteneryzacja & Docker, Docker Compose \\
Build & Gradle \\
CI/CD & GitHub Actions \\
\hline
\end{tabular}
\end{table}

\section{Podsumowanie}

Projekt systemu zarządzania projektami został zrealizowany zgodnie z założeniami, implementując:

\begin{itemize}
    \item Znormalizowany schemat bazy danych (3NF)
    \item Pełne więzy integralności i walidację danych
    \item Automatyczne tworzenie tabel przy pierwszym uruchomieniu
    \item System audytu z rejestracją czasu operacji
    \item Kompleksowe API REST z dokumentacją OpenAPI
    \item Bezpieczną autentykację opartą na JWT
    \item Powiadomienia i komunikację w czasie rzeczywistym
\end{itemize}

System jest gotowy do wdrożenia produkcyjnego i dalszego rozwoju.

% ============================================================================
% DODATKI
% ============================================================================
\appendix

\chapter{Pełny schemat SQL}

Pełny schemat bazy danych znajduje się w pliku:\\
\texttt{docs/database-schema.sql}

\chapter{Diagram przypadków użycia PlantUML}

Źródło diagramu w formacie PlantUML:\\
\texttt{docs/use-case-diagram.puml}

% ============================================================================
% BIBLIOGRAFIA
% ============================================================================
\begin{thebibliography}{9}

\bibitem{spring}
Spring Framework Documentation,
\url{https://docs.spring.io/spring-framework/docs/current/reference/html/}

\bibitem{hibernate}
Hibernate ORM Documentation,
\url{https://hibernate.org/orm/documentation/}

\bibitem{postgresql}
PostgreSQL Documentation,
\url{https://www.postgresql.org/docs/}

\bibitem{jwt}
JSON Web Tokens,
\url{https://jwt.io/introduction}

\bibitem{openapi}
OpenAPI Specification,
\url{https://swagger.io/specification/}

\end{thebibliography}

\end{document}
